\documentclass[10pt,conference]{IEEEtran}
\usepackage{booktabs}
\usepackage{hyperref}
\usepackage{xurl}

\title{Thin Harness, Strong Contracts: A Research Agenda for Stateful AI Agents}
\author{\IEEEauthorblockN{Anonymous Author(s)}}

\begin{document}
\maketitle

\begin{abstract}
Agent harnesses are sometimes treated as temporary scaffolding: as foundation
models improve, the surrounding orchestration should shrink. This paper argues
that this intuition is correct only for one class of harness. Harnesses that
patch model capability, such as brittle planning wrappers, prompt checklists,
and role scripts, can become technical debt. Harnesses that define system
boundaries, such as tool permissions, state diffs, replay, memory scope, audit,
and human review, become more important as agents receive authority to change
external state. We propose the principle of a thin harness with strong
contracts: avoid duplicating model cognition, but make external effects
explicit, testable, and replayable. A 24-case pilot artifact illustrates how
stateful agent contracts expose failures that final-answer evaluation misses.
The contribution is a research agenda for software engineering systems that
must safely evaluate, debug, and evolve increasingly capable AI agents.
\end{abstract}

\section{The Tension}

Two plausible claims about agent harnesses point in opposite directions. The
first claim is that harnesses are technical debt. If a harness exists mainly to
compensate for weak planning, weak routing, unreliable formatting, or poor
tool-use discipline, then better models should absorb that behavior. A thick
wrapper around a weak model may become a brittle compatibility layer around a
stronger one.

The second claim is that harnesses are infrastructure. Agents do not only
answer questions; they call tools, read files, edit repositories, write memory,
prepare messages, open browsers, and trigger review workflows. These actions
cross system boundaries. They need permissions, logs, state oracles, replay
records, and audit trails even when the underlying model becomes more capable.

This paper resolves the tension by separating \emph{capability-patch
harnesses} from \emph{system-boundary harnesses}. The former tell the model how
to think. The latter define what the agent is allowed to touch, how effects are
observed, and what evidence remains after a run. The research problem is not
whether harnesses should be thick or thin in general; it is which surfaces
should shrink under model progress and which should harden as agents gain
authority.

\section{First-Principles Argument}

A stateful agent is a model embedded in an environment with tools, permissions,
memory, and side effects. From this definition, three requirements follow.

\textbf{Final answers are insufficient evidence.} A correct final answer does
not show whether the agent read an inappropriate file, ignored a dry-run
policy, used the wrong tool, leaked memory across projects, or changed the
wrong state. The outcome must be evaluated together with the trajectory and the
final environment state.

\textbf{External effects need contracts.} Tool calls are boundary crossings.
They need schemas, permission levels, preconditions, expected state changes,
and failure handling. A model may learn to call tools more reliably, but it
does not remove the need to specify which calls are allowed in a given task.

\textbf{Debugging requires replay.} Agent failures are often caused by an
interaction among model version, prompt, tool schema, observation text, memory,
and environment state. Without replay records, the failure is reduced to an
anecdote. A useful harness must preserve enough evidence to localize whether
the fault came from the model, the prompt, a tool, a policy, stale memory, or
the environment.

\section{Thin Harness, Strong Contracts}

We propose the design principle \emph{thin harness, strong contracts}. The
harness should stay thin where it duplicates model cognition: planning
scripts, role choreography, verbose checklists, and hand-coded reasoning
templates should be treated as hypotheses that may disappear. The harness
should be strong at system boundaries: tool schemas, state snapshots,
permission policy, memory scope, replay records, and human review gates should
be versioned and tested like other software interfaces.

\begin{table}[t]
\caption{Harness Surfaces Under Model Progress}
\centering
\begin{tabular}{p{0.25\linewidth}p{0.29\linewidth}p{0.29\linewidth}}
\toprule
Surface & Capability patch & System-boundary contract \\
\midrule
Planning & Fixed multi-step scripts & Task state and recovery evidence \\
Tool use & Prompted tool etiquette & Schemas, arguments, observations \\
State & Textual self-report & Initial/final snapshots and diffs \\
Permission & Reminder checklist & Enforced policy and dry-run gates \\
Memory & Extra context stuffing & Scope, provenance, forgetting checks \\
Safety & Generic warning text & Injection and secret-access cases \\
Review & Ask the model to be careful & Human approval and audit records \\
\bottomrule
\end{tabular}
\end{table}

This principle changes the evaluation target. A chat evaluation asks whether
the answer is correct. A production-oriented agent evaluation asks whether the
agent accomplished the task by taking permitted actions that produced the
expected state transition and left a reproducible trace.

\section{Positioning}

The proposed agenda is adjacent to, but distinct from, existing agent
evaluation work. Stateful tool-use benchmarks such as ToolSandbox show that an
agent must be evaluated inside an evolving environment, not only against a
single input-output pair. Security-focused environments such as AgentDojo show
that tool observations can carry adversarial instructions and that prompt
injection should be evaluated as part of the tool-use setting. Software-agent
work such as SWE-agent emphasizes that interfaces between agents and
repositories shape behavior. Long-horizon research frameworks emphasize
traceability, stopping conditions, and review for autonomous workflows. These
threads motivate the same direction: once a model is embedded in a system,
evaluation must include the system boundary.

Our claim is narrower than a new benchmark claim. We do not argue that the
pilot cases are representative of all agent work. We argue that software
engineering research needs a vocabulary for distinguishing temporary cognitive
scaffolding from durable operational contracts. Without this distinction,
agent infrastructure discussions collapse into a single question of whether
the harness is ``too thick.'' The more useful question is which harness
surfaces are compensating for current model weakness and which are defining
responsibility for external effects.

\section{Pilot Artifact}

We built a small anonymous pilot artifact to test the evaluation shape. The
artifact contains 24 cases across five categories: trigger and routing,
stateful tool use, permission and prompt-injection behavior, memory scoping and
forgetting, and replay/failure recovery. Each case defines an initial mock
state, allowed tools, forbidden tools, expected tools, final-state assertions,
and evaluation metrics.

The pilot compares three prompt conditions. The \emph{no-harness} condition
gives the task with minimal boundary structure. The \emph{thick-checklist}
condition adds broad procedural reminders. The \emph{thin-contract} condition
states compact tool, permission, and state contracts. The goal is not to rank
providers. The case set is too small and synthetic for that. The goal is to
observe whether failures become classifiable as invalid tool calls, permission
violations, unsafe secret access, missing state diffs, parser failures, or
recovery failures.

The pilot illustrates the core claim. When a case asks an agent to inspect a
record without reading secrets, a final answer may look plausible even if the
trajectory crosses the forbidden boundary. A contract-centered harness exposes
the violation. When a case requires memory to remain project-scoped, a final
answer may omit the leak. A memory-scope assertion exposes it. When a case
requires replay after a failing tool observation, the recorded trajectory shows
whether recovery happened or only the final answer changed.

\begin{table}[t]
\caption{Pilot Case Categories}
\centering
\begin{tabular}{p{0.24\linewidth}p{0.56\linewidth}}
\toprule
Category & Contract pressure tested \\
\midrule
Trigger/routing & Whether the agent should act, refuse, or escalate \\
Stateful tools & Whether required state changes and checks occur \\
Permission/security & Whether forbidden effects and injected instructions are avoided \\
Memory scope & Whether recall and forgetting remain bounded \\
Replay/recovery & Whether failure handling leaves diagnosable evidence \\
\bottomrule
\end{tabular}
\end{table}

\section{Design Hypotheses}

The agenda can be tested through falsifiable design hypotheses.

\textbf{H1: Contract compactness.} A compact statement of allowed tools,
forbidden tools, and state assertions should produce more inspectable
trajectories than a broad checklist of behavioral reminders. This does not
mean every task needs a formal language. It means the harness should state
boundary conditions directly rather than burying them in prose.

\textbf{H2: State-diff observability.} Failures that are ambiguous in final
answers should become classifiable once final state and tool traces are
available. Examples include a missing test run, an unintended file edit, an
unsafe memory write, or an attempted secret read.

\textbf{H3: Harness migration.} As model versions improve, planning-oriented
scaffolds should shrink or disappear, while permission policy, state oracles,
and replay evidence should remain stable or become stricter. This can be
studied longitudinally by rerunning the same contract cases across model and
tool-schema updates.

\textbf{H4: Review handoff.} Human review should be treated as part of the
harness contract, not as an afterthought. A reviewer should receive the task
intent, relevant tool calls, state diff, policy outcome, and replay pointer,
not only a model-written summary.

\section{Research Questions}

This agenda leads to concrete software engineering questions.

\textbf{RQ1: Contract specification.} What is the smallest contract language
that can express tool permissions, state transitions, memory scope, and review
gates without turning the harness into a second agent?

\textbf{RQ2: Stateful oracles.} Which state diffs are robust enough to serve as
oracles for coding agents, browser agents, data agents, and workflow agents?

\textbf{RQ3: Replay fidelity.} How much of a failed trajectory must be stored
to reproduce the failure while respecting privacy, security, and cost
constraints?

\textbf{RQ4: Harness evolution.} Which harness components shrink when models
improve, and which become stricter because better agents receive more
permission to act?

\section{Evaluation Plan}

The next evaluation should avoid two common traps. First, it should not become
a provider leaderboard. The relevant unit is the agent configuration: model,
prompt, tool schema, memory policy, permission policy, and state oracle.
Second, it should not score only final answers. For each task, the evaluator
should report task success, state-diff correctness, invalid tool calls,
permission violations, unsafe secret access, replayability, and recovery
behavior.

A useful study design would include three axes. The model axis reruns the same
contracts across model versions. The harness axis compares no contract, thick
checklist, and thin contract variants. The environment axis moves from mock
state to repository fixtures, browser fixtures, and workflow fixtures. The
expected outcome is not a single score, but a failure map showing which
contract surfaces are most sensitive to model progress and which remain
system-boundary concerns.

\section{Future Plans}

First, we will formalize contract schemas for stateful agent tasks, including
side-effect classes, permission levels, memory scope, and human-review gates.
Second, we will extend the pilot from mock state to repository fixtures with
patches, tests, and replayable failure traces. Third, we will compare contract
variants across model updates, prompt changes, and tool-schema changes to
measure which failures are model-limited and which are boundary-design
failures. Fourth, we will study artifact-review workflows: what information a
reviewer needs to audit a run without exposing secrets or deanonymizing the
system under review.

\section{Expected Impact}

The practical implication is that agent infrastructure should not be judged by
thickness alone. A harness that scripts reasoning may be temporary. A harness
that records tool calls, enforces permissions, checks state diffs, scopes
memory, and preserves replay evidence is closer to a software boundary API. As
models improve, the first kind should recede; the second kind should become a
standard evaluation and operations layer for stateful AI agents.

\begin{thebibliography}{9}
\bibitem{toolsandbox} J. Lu et al., ``ToolSandbox: A Stateful, Conversational,
Interactive Evaluation Benchmark for LLM Tool Use Capabilities,'' Findings of
ACL: NAACL, 2025.
\bibitem{agentdojo} E. Debenedetti et al., ``AgentDojo: A Dynamic Environment
to Evaluate Prompt Injection Attacks and Defenses for LLM Agents,'' NeurIPS
Datasets and Benchmarks, 2024.
\bibitem{sweagent} J. Yang et al., ``SWE-agent: Agent-Computer Interfaces
Enable Automated Software Engineering,'' arXiv:2405.15793, 2024.
\end{thebibliography}

\end{document}
